\documentclass{article}
\usepackage[utf8]{inputenc}
\usepackage[T1]{fontenc}
\usepackage{hyperref}
\usepackage{graphicx}
\usepackage{amssymb}
\usepackage{amsmath}
%% \graphicspath{{../out/}}

\title{Colorful Components: MOP \\
\large Operations Research Complements
}
\author{Michael Chelli}
\date{\today}

\begin{document}
\raggedright
\maketitle
\newpage
\tableofcontents
\newpage
\section{Introduction}
Given a graph, where to each node is associated a color,
a colorful component of that graph is a connected subgraph such that no two nodes have the same color.

The "Minimum Orthogonal Partition" (MOP) problem consists in deleting the minimum number of edges from a graph
so that the remaining graph consists only of colorful components.

In (2024) is given an introduction to MOP, "Maximum Edges in transitive Closure" (MEC) and "Minimum Colorful Components" problems,
with their applications (mainly in biology and networks) and solutions based on branch and cut.

In particular, MOP was introduced in (2000), where an approximation algorithm was given, and an exact algorithm was first given in (2012)

\section{Mathematical formulation}
\subsection{Definitions}
In the MOP problem there is a graph
$G = (V,E,C)$ where $V$ is the set of nodes, $E$ the set of edges, and $C$ the set of colors.

$c : V \to C$ is a function giving the color of each node: for a node $u \in V$ its color is $c(u)$.

The problem consists in removing the minimum number of edges $E' \subset E$ such that in the resulting graph $G' = (V, E', C)$,
all connected components are colorful.

It is equivalent to find a way to partition the graph in $Q$ colorful subgraphs,
such that the sum of edges in each of them is maximized.
Note that with this formulation the subgraphs do not necessarly need to be connected.

\subsection{Integer Linear Programming formulation}
given the last definition, MOP can be formulized as follows:

\[
\begin{cases}
	\max \sum_{q \in Q}{\sum_{ei = 1}^{|E|}{e_{q,ei}}} \quad \text{s.t.}\\
	x \in \{0,1\} ^{ Q \times |V| }\\
	e \in [0,1]^{ Q \times |E| }\\
	\text{(1) } \forall i \in V & \sum_{q=1}^{Q}{x_{q,i}}\le 1 \\
	\text{(2) } \forall q \in Q \forall k \in C & \sum_{i \in V : c(i) = k}{ x_{q,i} } \le 1 \\
	\text{(3) } \forall q \in Q \forall (ei,u,v) \in E & e_{q,ei} \le x_{p,u} \land e_{q,ei} \le x_{q,v}\\
\end{cases}
\]

where in (3) $ei$ is an index associated to the edges ranging from $1$ to $|E|$.
And $Q$ is a bound on the number of subgraphs for an optimal solution, since the exact number is hard to know.

\subsection{Bound on Q}
A simple upper bound to $Q$ would be $N$, in (2024) better bounds are shown for when the subgraphs need to be connected.

If, like in the above formulation, the subgraphs need not be connected, we can do better:\\
define $nc : C \to N$ where $nc(k) = | u \in V : c(u) = k |$.\\
Call a partitioning to be minimal if all pairs of subgraphs have at least one color in common.\\
There is always an optimal solution that is minimal: if it is not you can merge two subgraphs that have no color in common.\\
Consider the graph of a partition, where each node is a subgraph and there is an edge between two subgraphs if they have at least one color in common.\\
In a minimal solution, such graph is a clique, so it has $q = \frac{s \cdot (s-1)}{2}$ edges if $s$ is the number of subgraphs, $s = \sqrt{2q + \frac{1}{4}} + \frac{1}{2}$ .\\
Each color $k$ in the original graph can make at most $\frac{nc(k) \cdot (nc(k)-1)}{2}$ edges.\\
So there are at most $K = \sum_{k \in C}{\frac{nc(k) \cdot (nc(k)-1)}{2}}$ edges.\\
thus a minimal solution always has at most $Q = \lfloor \sqrt{2K + \frac{1}{4}} + \frac{1}{2} \rfloor$ subgraphs.

\section{Branch and price algorithm}
\subsection{Master and Pricing problems}
In the ILP formulation, (1) are linking constraints, we can separate these from the others obtaining a Dantzig-Wolfe decomposition with identical subproblems.

We convexify constraints (2) and (3), and do column generation, where constrains in (1) remain in the master problem and the other go in the pricing subproblem.

The master problem is:
\[
\begin{cases}
	\max \sum_{u=1}^{p} \sum_{q=1}^{Q} \sum_{ei = 1}^{|E|} e^{(u)}_{q, ei} \cdot \theta_u \quad \text{s.t.}\\
	\forall i \in V \quad \sum_{u=1}^{p} \sum_{q=1}^{Q} x_{p,i}^{(u)} \cdot \theta_u \le 1\\
\end{cases}
\]
and the pricing subproblem is:
\[
\begin{cases}
	\max \sum_{i \in V} \lambda_i + Q(\sum_{ei = 1}^{|E|}{e_{p,ei}} - \sum_{i \in V}{\lambda_i x_i}) \quad \text{s.t.}\\
	\forall q \in Q \forall k \in C \quad \sum_{i \in V : c(i) = k} x_{q,i} \le 1\\
	\forall q \in Q \forall (ei,u,v) \in E \quad e_{q,ei} \le x_{q,u} \land e_{q,ei} \le x_v\\
\end{cases}
\]
This is equivalent to a lagrangean relaxation of the ILP formulation, where (1) is lagrangeanized.
So whenever we solve a pricing subproblem optimally, we also obtain a dual bound for free.

The $Q$ dimension is now superfluous and the problem can be rewritten as:
\[
\begin{cases}
	\max \sum_{i \in V} \lambda_i + Q(\sum_{ei = 1}^{|E|}{e_{p,ei}} - \sum_{i \in V}{\lambda_i x_i}) \quad \text{s.t.}\\
	\forall k \in C \quad \sum_{i \in V : c(i) = k} x_{i} \le 1\\
	\forall (ei,u,v) \in E \quad e_{q,ei} \le x_{u} \land e_{ei} \le x_v\\
\end{cases}
\]

We can get rid of it in the master problem as well by having $\theta$ sum to $Q$ instead of $1$:
\[
\begin{cases}
	\max c \cdot \theta \quad \text{s.t.}\\
	\forall i \in V \sum_{u=1}^{p} \theta_u \cdot x_i^u \le 1\\
\end{cases}
\]

And the pricing subproblem's objective will change to:
$\max \sum_{i \in V} \lambda_i + \sum_{ei = 1}^{|E|}{e_{ei}} - \sum_{i \in V}{\lambda_i x_i}$

And any optimal solution for this will also be an optimal solution for the lagrangean relaxation.
\subsection{Braching strategy}
I branch on edges, like is done in (2012).
On one side of the branch, the chosen edge will be deleted, on the other side it will be merged.
This means that in the pricing problem, the edges will now be weighted, and each node can have multiple colors associated to it.

Branching is done on the edge whose $e_{ei}$ is closest to $0.5$ in the RLMP solution.
I first explore the branch where the edge is merged, this was somewhat of an arbitrary decision.

After solving the LMP, I solve the Restricted Master Problem with ILP, which usually gives good primal bounds,
often good enough to close the BnP on the first node.

I always keep all generated columns, and feed all compatible ones into each instance of the master problem.

\subsection{Hotstart}
Columns were hotstarted with heuristic solutions.

Two different heuristic approaches were considered: Karger's heuristic(ref) and Merge heuristic(ref),
the first of which consistently gave slighthly better results.
Only the one set of best columns were used to hotstart the first node of the BnP.

\subsection{Stabilization}
I used single parameter stabilization(ref) %% TODO: write more

\section{Pricing subproblem}
The pricing subproblem consists in finding the colorful subgraph with the biggest edge weight sum minus node costs.
Where each node has multiple colors and a real cost, and each edge has an integer weight.

For each edge $(u,v)$, define $w(u,v)$ (or $w(ei)$) as its weight, or $0$ if there is no such edge.

For each node $u$, call $\lambda_u$ its cost.

\subsection{Preprocessing}
For simplicity, consider each node to have a single color during the preprocessing.

A node $u$ is useless if $-\lambda_u + \sum_{k \in C} \max_{v : c(v)=k} w(u,v) < 0$

An edge $(u,v)$ is useless if $-\lambda_u -\lambda_v + \sum_{k \in C \setminus \{c(u),c(v)\}} \max_{i : c(i)=k} w(u,i)+w(v,i) < 0$

We can iteratively delete useless nodes and edges to obtain a smaller problem.

\subsection{ILP formulation}
\[
\begin{cases}
	\max \sum_{ei = 1}^{|E|}{e_{ei} \cdot w(ei)} - \sum_{u \in V} \lambda_u \quad \text{s.t.}\\
	x \in \{0,1\}^{|V|}\\
	e \in [0,1]^{|E|}\\
	\forall u \in V \forall v \in V : c(v) = c(u), v \neq u \quad x_u + x_v \le 1\\
	\text{(3) }\forall (ei,u,v) \in E \quad e_{ei} \le x_u \land e_{ei} \le x_v\\
\end{cases}
\]

\subsection{Branch and Bound solution}
I only consider connected components in this solution, which has the advantage that it restricts the search, especially for sparse graphs,
and the disandavtage that we can no longer use the tighter bound we obtained on $Q$.

I branch on the nodes and keep a connected set of "taken" nodes $U$.

Call $C'$ the set of colors that do not appear in $U$.

Any node further than $C'$ from $U$ can be discarded, as to reach it you would have to take the same color twice at least once.

Any node $u$ with $-\lambda_u + \sum_{v \in U}{w(u,v)} + \sum_{k \in C'}{\max_{v : c(v)=k} w(u,v)} < 0$ is useless and can be discarded.

I split each edge's weight across its two nodes, this is a surrogate relaxation for each edge of constraints (3) in the previous formulation.\\

If each edge $w(u,v)$ was split with $0 \le \alpha \le 1$ value goint towards $u$ and $1 - \alpha$ going towards $v$,
call $w^{(u)}(u,v) = \alpha \cdot w(u,v)$ and $w^{(v)}(u,v) = (1-\alpha) \cdot w(u,v)$.

The value of a node $u$ will now be: $\sigma(u) = -\lambda_u + \sum_{v \in U}{w(u,v)} + \sum_{k \in C'}{\max_{v : c(v)=k} w^{(u)}(u,v)}$.\\
And $\sum_{k \in C'}{\max_{u : c(u) = k}{ \sigma(u) }}$ is a valid dual bound.

In a couple of places (where I take the max across colors) I considered each node to have a single color, like in the preprocessing.
This could lead to looser bounds and it might be better to split each node along its colors.

It should not be hard to find good values for $\alpha$ with a local search, but I just used an (quite arbitrary) heuristic:
considering $w(u,v)$, its $\alpha$ is $\frac{4 + \lambda_u}{8 + \lambda_u + \lambda_v}$.

\subsection{Heuristic pricing}
I considered several heuristic approaches to pricing.

In the end I kept a beam-search based approach,
where each entry is a connected subgraph and I incrementally add the next node that increases the score the most.

When doing heuristic pricing, the lagrangean dual value cannot be updated,
and the columns are not as good as with the exact approach,
but it seems to speed up solving the master problem on average.


\section{Experiments}
All runs were executed on a "11th Gen Intel(R) Core(TM) i5-1135G7 @ 2.40GHz" CPU, with a time limit of 180s per graph.

\subsection{Dataset 1}

This first dataset is from (2012)%%, it has been made from TODO.

It contains 135 graphs, each with several connected components. When there are multiple components, each one is solved with a separate BnP.

\begin{center}
\begin{tabular}{ |c|c|c|c| } 
 \hline
 & bnp & bnp\_hp & bnp\_nohot \\
 \hline
 solved & 112 & 113 & 110 \\
 \hline
 avg time & 5.3s & 4.2s & 7.2s \\
 \hline
 avg pricing bnb nodes & 374255 & 235378 & 518371 \\
 \hline
\end{tabular}
\end{center}

\subsection{Dataset 2}
2024 data

\subsection{Analysis}
%% analyze changing N, density, C

%% Aknowledgments (highs, graphs), link repo
\section{Aknowledgments}

For solving all LP and ILP instances, HiGHS(ref) was used.

Graphs for dataset 1 and 2 have been taken from (2012) and (2024) respectively.

The code is available at(), and can be compiled with make, HiGHS needs to be installed.

\end{document}